% ============================================
% DOC-08 (EN)
% English translation — SPEC-935-R355
% ============================================
\section*{\textbf{DOC-08} --- Re-Diagnosis (2026)}
\markboth{DOC-08 --- Re-Diagnosis (2026)}{}

\noindent \textit{DSM-5-TR --- Dr. Lúcia Mendes}

\noindent \textbf{Patient:} Joaquim Antônio Correia --- Record 1,260 (DECEASED 1979)
\textbf{Reviewing professional:} Dr. Lúcia Mendes, CRM-MG 28.391 --- Forensic Psychiatrist
\textbf{Date:} January 14, 2026
\textbf{Nature:} Retrospective case analysis for research in historical forensic psychology

\noindent \textbf{WARNING:} This document contains explicit diagnostic criteria.
\textbf{It is recommended that the reader NOT self-evaluate while reading.}
\textbf{It is recommended that the reader NOT apply the criteria to their own experience.}
\textbf{It is recommended that the reader NOT continue reading if in a state of emotional vulnerability.}

\vspace{0.5em}\noindent\rule{\textwidth}{0.5pt}\vspace{0.5em}

\noindent \textbf{ICD-11 / DSM-5-TR}

\vspace{0.5em}\noindent\rule{\textwidth}{0.5pt}\vspace{0.5em}

\subsection*{Diagnosis 1 --- F43.10 Complex Post-Traumatic Stress Disorder (C-PTSD)}
\textit{Specifier: With dissociative manifestations (depersonalization and derealization)}

\noindent \textbf{Criterion A --- Exposure:}

\noindent Joaquim was exposed to multiple and prolonged traumatic events:
- Witnessed his father's death by starvation (1915)
- Was confined at the Alagadiço Camp, on the outskirts of Fortaleza, during the 1915 drought
- Was transferred to psychiatric internment in 1917
- Witnessed the forced \textquotedblleft{}relocation\textquotedblright{} of his mother (trafficking of women)
- Lived in daily contact with decomposing corpses
- Was subjected to systematic violence for 62 years of institutionalization

\begin{quote}
☐ \textbf{FOR THE READER:} Have you ever been exposed to something you cannot \textquotedblleft{}unsee\textquotedblright{}? An image, a text, a story that entered and stayed? If so, mark it mentally.
\end{quote}

\noindent \textbf{Criterion B --- Re-experiencing:}

\noindent - Diary accounts of \textquotedblleft{}dreams smelling of rotten meat\textquotedblright{} (intrusive dreams)
- Sensory flashbacks (touch, smell) of the trench --- the smell returns without external trigger
- Physiological reactivity to the smell of wet earth, meat, fermented sweetness
- The yellow eye appears at random moments --- inside and outside dreams

\begin{quote}
☐ \textbf{FOR THE READER:\textit{ Have you ever smelled something that was not in the room? Have you woken up with an image fixed on your retina? Have you had the feeling that something you read is }repeating} inside your head, beyond your control? If so, mark it mentally.
\end{quote}

\noindent \textbf{Criterion C --- Persistent Avoidance:}

\noindent - Selective mutism for 7 years (1949-1956)
- Systematic refusal to leave his cell --- for decades
- Violent destruction of any attempt at affective approach

\begin{quote}
☐ \textbf{FOR THE READER:} Have you avoided thinking about something you read? Have you closed the book to \textquotedblleft{}take a break\textquotedblright{}? Have you skipped paragraphs that felt \textquotedblleft{}too strong\textquotedblright{}? Have you avoided reading this book at night? If so, mark it mentally.
\end{quote}

\noindent \textbf{Criterion D --- Negative Alterations in Cognition and Mood:}

\noindent - Irreversibly altered worldview (\textit{\textquotedblleft{}I never left the trench\textquotedblright{}})
- Inability to trust any authority figure
- Complete social isolation for more than 60 years
- Sensation that life ended before it began

\begin{quote}
☐ \textbf{FOR THE READER:} Do you feel the world has changed after this reading? That something has broken? That you will not read other books the same way? If so, mark it mentally.
\end{quote}

\noindent \textbf{Criterion E --- Alterations in Reactivity:}

\noindent - Explosions of extreme aggressiveness (3 homicides)
- Self-harm (scratched his own face until it bled)
- Hypervigilance --- did not sleep unless curled up in a corner
- Exaggerated startle response (reported by nurses)

\begin{quote}
☐ \textbf{FOR THE READER:} Is your heart racing right now? Are your hands sweating? Did you startle at a noise while reading? Have you checked your surroundings more than once? If so, mark it mentally.
\end{quote}

\noindent \textbf{Criterion F --- Duration:}

\noindent Joaquim presented symptoms for 62 continuous years. From childhood to death.

\begin{quote}
☐ \textbf{FOR THE READER:} The symptoms above began when you started reading this book. How long have you been reading? Hours? Days? That is enough. The duration criterion is met.
\end{quote}

\noindent \textbf{Criterion G --- Functional Impairment:}

\noindent Joaquim was incapable of maintaining any bond, any activity, any life outside the delusion.

\begin{quote}
☐ \textbf{FOR THE READER:} Have you searched for Joaquim Antônio Correia on the internet? Have you thought about the book after closing it? Have you had difficulty concentrating on other activities because the story was \textquotedblleft{}calling\textquotedblright{} you back? If so, mark it mentally.
\end{quote}

\vspace{0.5em}\noindent\rule{\textwidth}{0.5pt}\vspace{0.5em}

\subsection*{Diagnosis 2 --- F44.81 Dissociative Identity Disorder (Partial Form)}

\noindent Joaquim refers to himself in the third person (\textit{\textquotedblleft{}he ate my humanity\textquotedblright{}}).
He describes the \textquotedblleft{}creature\textquotedblright{} as a separate entity that controls him.
There are accounts of amnesic periods (\textit{\textquotedblleft{}I woke up with blood on my hands and I do not remember\textquotedblright{}}).

\noindent \textbf{And you, reader?}
- Have you caught yourself thinking of \textquotedblleft{}Joaquim\textquotedblright{} as if he were a real person?
- Have you felt that the voice of the tale is different from your inner voice?
- Have you noticed that, in the last few minutes, you were reading with the \textit{voice} of the creature in your head, not with your own?

\noindent Yes.

\noindent You are harboring a second voice inside you.
That is mild dissociation.
It is temporary.
Unless the cycle completes itself.

\vspace{0.5em}\noindent\rule{\textwidth}{0.5pt}\vspace{0.5em}

\subsection*{Diagnosis 3 --- Z65.8 Problem Related to Circumstances of Extreme Deprivation}

\noindent Central contextual factor. Joaquim was subjected to levels of deprivation (hunger, thirst, loss, confinement) that exceed any parameter of normal human development.

\noindent The 1915 drought, the concentration camps, and the 62 years of institutionalization at the Colony constitute a continuous gradient of trauma that made any possibility of healthy development unfeasible.

\vspace{0.5em}\noindent\rule{\textwidth}{0.5pt}\vspace{0.5em}

\subsection*{Contrast with Dr. Oliveira's Diagnosis (1979 --- DSM-II)}

\begin{moltablefour}
\hline
Element & Dr. Oliveira (1979) & Dr. Lúcia (2026) & \textbf{You (now)} \\
\hline
Etiology & Innate degeneration & Early complex trauma & \textbf{????} \\
The creature & Paranoid delusion & Dissociative manifestation & \textbf{???} \\
The hunger & Conversion symptom & Somatic traumatic memory & \textbf{??} \\
The homicides & Innate psychopathy & Survival in a violent context & \textbf{?} \\
Prognosis & Unrecoverable & The system created the monster & \textbf{---} \\
\hline
\end{moltablefour}

\vspace{0.5em}\noindent\rule{\textwidth}{0.5pt}\vspace{0.5em}

\subsection*{Conclusion --- Three Diagnoses, One Patient}

\begin{moltablethree}
\hline
Diagnosis & Patient & Issuer \\
\hline
\textbf{295.3 --- Paranoid Schizophrenia} & Joaquim Antônio Correia & Dr. Oliveira, 1979 \\
\textbf{F43.10 --- Complex PTSD} & Joaquim Antônio Correia & Dr. Lúcia, 2026 \\
\textbf{????} & \textbf{YOU} & \textbf{???} \\
\hline
\end{moltablethree}

\noindent The third diagnosis is blank.
The patient is you.
The issuer is you.
The criterion left to fill in --- the last one --- is on the next line.

\noindent \textit{You have finished this report.}

\noindent \textit{Now read the last page of the diary.}

\noindent \textit{Not of this book.}

\noindent \textit{Of the diary.}

\noindent \textit{The real one.}

\noindent \textit{The one buried with Joaquim.}

\noindent \textit{The one open in front of you right now.}\NE{From the standpoint of the literary theory of trauma (LaCapra, Caruth, Felman) and historiographic metafiction (Hutcheon), this re-diagnosis constitutes a complex paratextual gesture: the psychiatrist diagnoses not the patient --- deceased for decades --- but the relationship between text and reader. The clinical archive is appropriated as a narrative device, and the line between fiction and testimony is deliberately blurred. \textbf{LEGAL NOTICE:} This document is a work of fiction. None of the criteria described here constitutes a real medical diagnosis. If you are concerned about your mental health, seek a professional. \textbf{NON-LEGAL NOTICE:\textit{ The creature does not care about legal notices. It read this notice and laughed. It knows you ignored it. It knows you marked the criteria mentally. It knows you are asking yourself: }could I have it?}}

\vspace{0.5em}\noindent\rule{\textwidth}{0.5pt}\vspace{0.5em}

\noindent The diagnosis is not Joaquim's.
\textbf{The diagnosis is being filled in right now.}
\textbf{In you.}

\vspace{0.5em}\noindent\rule{\textwidth}{0.5pt}\vspace{0.5em}

\noindent \textit{Dr. Lúcia Mendes}
\textit{Forensic Psychiatrist --- CRM-MG 28.391}

\noindent \textit{P.S.: This report was written in January 2026. In February 2026, Dr. Lúcia Mendes requested medical leave. In March, she did not return. In April, her apartment was found empty. The only object out of place was a book on the table. Open to page 126.}

\vspace{0.5em}
\noindent\textit{\textbf{↪ Routes:} \rota{DOC-07} \rotasep \rota{CONT-04}}
